\section{Problem Formulation \& 
Prerequisite} 
\label{sec:problem-formulation-and-approach}
\subsection{Problem formulation}
The deployment of a real-time electricity price forecasting
system requires choices about three intertwined questions:
which model architecture to use, which data sources to acquire
and maintain, and over what horizon predictions remain useful.
While prior work ~\cite{energyPricePredictionMLNeighbors, 
energyPricePredictionUK, energyPricePredictionNovelDeepLearning} has addressed 
each question in isolation, under specialised 
settings, the practical 
interactions between these choices remain under-explored. A company
deploying such a system must weigh the recurring cost of
data acquisition against the marginal predictive accuracy each
data source provides, often without clear guidance on whether
comprehensive feature collection is justified.

This paper investigates how the predictive performance of
common machine learning architectures responds to 
systematic
variation in three settings: feature composition, forecasting
horizon, and prediction target. Formally, the problem is
structured as follows:

Given a set of architectures $M$ = RF, XGB, LGBM, CB,
LSTM, GRU, Transformer, Gluon, a set of forecast horizons
$H$ = 0, 24, 48, 72, 96, 120, 144, 168 hours, a set of feature
subsets $F$ (13 progressively expanding feature constellations),
and two prediction targets $T$ = Price, Delta, we evaluate the
performance of every time-feasible combination $(m, h, f, t) \in
M \times H \times F \times T$ using walk-forward validation over 11 years
of Danish electricity prices.

The central research question is: How much, and which,
data does a forecasting system actually need, to perform
competitively on short-term electricity price prediction in
Denmark?
This is investigated through two sub-questions:
\begin{enumerate}
	\item Architectural sensitivity: Do different architecture families respond 
	differently to changes in feature availability
	and forecasting horizon?
	\item Feature contribution: Which feature categories are most
	impactful to predictive accuracy, and at what point does
	adding more data stop yielding improvements?
\end{enumerate}
Performance is measured using MAE in EUR/MWh, chosen
for its interpretability and robustness to the frequent positive and negative 
price spikes characteristic of the Danish
electricity market. A naive persistence baseline 
serves
as the non-ML reference point. Supplementary experiments
address potential methodological pitfalls, including default
hyperparameter assumptions, loss function alignment, weather
data source addition, and cross-region generalisation.
This study deliberately does not pursue peak performance
through exhaustive hyperparameter tuning or architectural
specialisation. The contribution lies instead in the comparative behaviour of 
default-configured models across a wide
experimental grid, yielding practical guidance for the data-versus-performance 
trade-off facing real-world deployments.


\subsection{Forecasting approach selection}
As previously mentioned, the forecasting approaches chosen for predicting the 
electricity prices in Denmark are:
\begin{multicols}{2}
	\begin{enumerate}
		\item \gls{rf}
		\item \gls{xgb}
		\item \gls{lgbm}
		\item \gls{cb}
		\item \gls{lstm}
		\item \gls{gru}
		\item Transformer
		\item \gls{autogluon}
	\end{enumerate}
\end{multicols}
These approaches can be divided into three general 
groups: The tree-based models (RF, XGB, LGBM, and 
CB), the Neural network-based models (LSTM, GRU, 
and Transformer), and Gluon. The tree based 
models share several similarities, such as their 
prediction approach being based on decision trees, 
as well as their fast training times. The neural network-based models were 
selected due to their inherent memory capabilities (e.g., hidden states and 
attention mechanisms), which conceptually allow them to capture complex, 
long-term temporal dependencies in sequential data. However, due to their 
architectural complexity and the computational expense of sequential 
processing, they exhibit significantly longer training times compared to 
tree-based algorithms \cite[Sec. 2]{neuralNetworksMemory}. 
Gluon gets its own category, due to its 
inherent differences to the other models, as it is 
not a model itself, but an ensemble model creator. 
Further, it utilises decision trees, along with 
other base-models, for its forecasting ensemble, 
but also exhibits longer training-times similarly 
to the neural network-based models.\\

To maintain a standardised baseline across the $1,052$ experiments performed 
for this paper, as detailed in \autoref{sec:experiments}, all models were 
implemented using their default configurations, with no changes to 
hyperparameters or loss-functions. While it is acknowledged that neural 
networks are highly sensitive to hyperparameter tuning and may underperform in 
their default states compared to tree-based algorithms, performing exhaustive 
hyperparameter optimisation for every data 
constellation and horizon was 
computationally prohibitive. Consequently, the purpose of this paper is not 
to compare peak model performance against one another. The goal is to 
evaluate their relative performance by examining how the accuracy of 
different algorithmic families scales or degrades in response to varying data 
availability and increasing forecasting horizons. Aligning with this approach, 
AutoGluon was constrained to running on its fastest quality setting, trading 
accuracy for training speed, in order to ensure its training time would not 
increase prohibitively while performing 
walk-forward validation. To assess the impact of 
these simplifying choices, a supplementary set of 
experiments was conducted. First, the 
best-performing neural network model was retrained 
using MAE loss instead of the default MSE loss, to 
evaluate whether loss function alignment with the 
evaluation metric would improve performance. 
Second, hyperparameter tuning via Optuna was 
performed on the best, mean, and worst tree-based 
feature sets at the 24h horizon to quantify the 
improvement margin available through optimisation. 
Third, AutoGluon was retrained at the 
medium\_quality preset to determine whether the 
computational constraints of the fastest setting 
had materially limited its accuracy. Results from 
these supplementary experiments are discussed in 
\autoref{sec:results}.

\subsection{Prerequisite}\label{sec:prereq}
As previously mentioned, the Danish electricity 
market is special in various 
ways. These circumstances will be discussed here.
\subsubsection{Daily market fixing}
The most important thing to note about the Danish energy market is that the 
hourly energy prices are fixed at noon via the day-ahead-market. At noon, Nord 
Pool stops accepting bids and offers, processes all agreements and trades, and 
then publishes the price based on this somewhere between 30-90 minutes later. 
After this, energy producers, consumers, and 
traders, may utilise the 
intra-day-market to minimise the fines and 
penalties that come with producing 
too much or too little depending on the actual demand. For the purposes of this 
paper, we only consider the day-ahead-market.
\subsubsection{Regions}
The Danish energy grid is divided into two distinct regions/markets: DK1 and 
DK2, west and east of the Great Belt respectively. These regions trade with 
neighbouring countries independently, as well as trade with each other. For the 
purposes of this paper, it simply allows us to distinguish which weather 
stations' measurements should be used to predict the energy prices of the 
appropriate region.
\subsubsection{Time format}
As of October 2025, Nord Pool changed their reporting format for energy 
prices from hourly to 15-minute increments. As this paper uses data that spans 
both reporting formats, we simply discard data that is not exactly on the hour, 
as interpolating 15-minute increments for almost 10 years of data seemed 
unnecessary compared to the chosen approach. 